\documentclass[journal]{IEEEtran}
\usepackage{cite}
\usepackage{amsfonts}
\usepackage{amssymb}
\usepackage{amsmath}
\usepackage{amsthm}
\usepackage{url}
\usepackage{enumitem}
\usepackage{array}
\usepackage{booktabs}
\usepackage[utf8]{inputenc}
\usepackage[english]{babel}

\newtheorem{definition}{Definition}
\newtheorem{assumption}{Assumption}
\newtheorem{lemma}{Lemma}
\newtheorem{corollary}{Corollary}
\newtheorem{law}{Law}
\newtheorem{theorem}{Theorem}
\newtheorem{proposition}{Proposition}

\newcommand{\Volatile}{\mathsf{Vol}}
\newcommand{\Cached}{\mathsf{Cch}}
\newcommand{\Durable}{\mathsf{Dur}}
\newcommand{\Rebuild}{\mathsf{Reb}}
\newcommand{\Retained}{\mathsf{Ret}}
\newcommand{\Ok}{\mathsf{Ok}}
\newcommand{\Err}{\mathsf{Err}}
\newcommand{\Exhausted}{\mathsf{Exhausted}}
\newcommand{\El}{\mathsf{E}}
\newcommand{\Rm}{\mathsf{R}}
\newcommand{\dg}{h}
\newcommand{\ret}{\operatorname{retire}}
\newcommand{\live}{\operatorname{live}}
\newcommand{\stab}{\operatorname{stab}}
\newcommand{\enc}{\operatorname{enc}}
\newcommand{\Tmax}{T_{\max}}
\newcommand{\Dmax}{D_{\max}}
\newcommand{\Pmax}{P_{\max}}
\newcommand{\issue}[1]{\#{#1}}

\author{Ryan~J.~Kung%
\thanks{Ryan J. Kung can be reached at ryankung@ieee.org.
Manuscript created October 24, 2026.  This document specifies when a removal of a data-topic
element of Rings is stable, what may drop it, how large it may grow, how the fetch cache is
linearized, and what a registry heartbeat writes, before 1.0 (\texttt{rings-core}, modules
\texttt{dht::entry}, \texttt{dht::chord::storage} and \texttt{storage}; \texttt{rings-node},
module \texttt{registration}).  It is the law for issues \issue{918}, \issue{920},
\issue{921} and \issue{922} of \texttt{RingsNetwork/rings}, and extends the durable-records
law~\cite{DurableRecords}.}}
\title{\huge Removal Stability:\\
Eviction, Tombstone Bounds and Heartbeats\\
for the Data Topics of Rings}
\date{}

\IEEEaftertitletext{%
\vspace{0.5\baselineskip}%
\noindent\hfill
\parbox{0.90\textwidth}{%
{\centering\bfseries Abstract\par}%
\vspace{0.4\baselineskip}%
\noindent
A Rings data topic is a replicated last-writer-wins element set whose removes are collected
at a clock horizon (\issue{872}, \issue{913}).  Four open defects share one cause: nothing
states which actions may drop a remove before it is stable, nor how many removes a carrier
may hold.  Storage eviction drops removes and resurrects removed elements (\issue{920}); a
sync peer can grow a carrier's removes past one transport message (\issue{921}); the fetch
cache was specified without a linearization (\issue{922}); and registry heartbeats emit one
remove each, about 25\,KB per registrant (\issue{918}).  We define stability, state one
survival law for removes, and prove no resurrection under it.  We factor a carrier into an
element side and a remove side, give each its point of the durable-records authority
lattice, and so confine eviction to elements.  We bound the remove side by a count enforced
by refusal, prove the carrier fits one message, and show that a forged sync delta simulates
a writer.  We give the cache's linearization, and replace the heartbeat by the re-append of
a signed lease, which removes nothing and reads nothing.
}\hfill\null\par
\vspace{0.5\baselineskip}}

\begin{document}
\maketitle
\emergencystretch=1.5em

\section{Problem}\label{sec:problem}

\begin{description}[leftmargin=1.2em,style=nextline]
\item[\issue{920}] The byte budget of the data store evicts the least recently written
carrier, whatever it holds.  An evicted remove lets a stale replica's sync re-add the element
it covered.
\item[\issue{921}] A data carrier's removes have no count or byte bound, and a sync delta
may carry any number of them.  The assertion that a full carrier fits one transport message
counts elements only.
\item[\issue{922}] A cache read that retires an entry can overwrite or delete a value that a
concurrent cache put just stored.
\item[\issue{918}] A registry descriptor's bytes change on every heartbeat, so every
heartbeat removes the previous descriptor.  Removes are retained for $H+\sigma$, which leaves
about $200$ removes per registrant, and every registrant reads all of them on every heartbeat.
\end{description}
\issue{920} and \issue{921} concern the survival and the size of a remove.  \issue{918} is
the dominant producer of removes.  \issue{922} is the one store that holds carriers and is
not replicated.  Section~\ref{sec:model} fixes the model, Section~\ref{sec:stable} the
survival law, and Sections~\ref{sec:evict}--\ref{sec:heartbeat} one law per issue.
Section~\ref{sec:plan} orders the implementation.

\section{Model}\label{sec:model}

\begin{definition}[Carrier]\label{def:carrier}
A data carrier is $x=(V,D,\rho,\beta)$: adds $V:\mathrm{Bytes}\rightharpoonup\mathrm{Dot}$,
removes $D:\mathrm{Digest}\rightharpoonup\mathrm{Dot}$, a register
$\rho\in\mathrm{Ver}_\bot$ (the overwrite floor) and a bound $\beta\in\mathbb{N}_\bot$.  A dot
leads with its hybrid time $\tau(d)$.  The remove $(e,r)\in D$ \emph{covers} the add
$(v,d)\in V$ iff $\dg(v)=e\wedge d\le r$, with $\dg$ Keccak-256.  $x$ is \emph{normalized}
iff no add is covered or below $\rho$ and no remove is below $\rho$.  The join $\sqcup$ is
the pointwise $\max$ followed by normalization (\texttt{crdt.rs}).  It is a semilattice on
normalized carriers~\cite{Shapiro2011}.
\end{definition}

\begin{definition}[Horizon, retirement, liveness]\label{def:retire}
With $H$ the element horizon ($6\cdot10^6$\,ms) and $\sigma$ the skew tolerance
($3000$\,ms), as in \issue{913}:
\begin{align*}
\ret_t(x)&=x\setminus\{(v,d):t\ge\tau(d)+H\}\\
&\phantom{{}=x}\setminus\{(e,r):t\ge\tau(r)+H+\sigma\}\\
&\phantom{{}=x}\setminus\{\rho:t\ge\tau(\rho)+H+\sigma\},\\
s(x)&=\max\{\tau(r):r\in\operatorname{ran}D\cup\{\rho\}\}+H+\sigma,\\
\live_t(x)&\iff\beta\ne\bot\wedge(t<\beta\vee t<s(x)).
\end{align*}
$\ret_t$ is a join homomorphism and $\ret_t\circ\ret_s=\ret_{\max(s,t)}$.
\end{definition}

\begin{assumption}[Clocks]\label{as:clock}
Any two node clocks differ by at most $\sigma$.  Admission at clock $c$ refuses a delta that
carries a version with $\tau>c+\sigma$.
\end{assumption}

\textbf{Actors.}  A \emph{writer} is any node: a data topic has no write authority, so any
node may $\mathrm{Extend}$, $\mathrm{Tombstone}$ or $\mathrm{Overwrite}$ any topic.  A
\emph{syncer} is any node that sends a delta (a hand-off, a repair, a lookup reply), which
the receiver admits by the admission law alone.  An \emph{owner} holds a replica of a
placement in its data store.  The fetch cache is a \emph{mirror}: it is not a replica, it
never acknowledges, and it spreads only by read-repair.

\textbf{Program actions and loss events.}  A program action is a transition of the store
that the program performs: a join, an operation, a retirement, an ack-gated deletion,
eviction, a cap, or a refusal.  A loss event is what the platform or the operator does
(Definition~4 of~\cite{DurableRecords}).  The distinction matters because a writer controls
program actions through the writes it sends, but it cannot cause a loss event.

\section{Stability}\label{sec:stable}

\begin{definition}[Stable]\label{def:stable}
A remove $(e,r)$ is \emph{stable} at clock $t$ iff $t\ge\tau(r)+H+\sigma$, and likewise the
register $\rho$ iff $t\ge\tau(\rho)+H+\sigma$.  A carrier's remove side is stable at $t$ iff
$t\ge s(x)$.
\end{definition}

\begin{lemma}[A stable remove covers only retired adds]\label{lem:covers}
If $(e,r)$ is stable at the clock $t_p$ of some node, then every add $(v,d)$ it covers is
retired at the clock $t_q$ of every node at the same instant.
\end{lemma}
\begin{proof}
$t_q\ge t_p-\sigma\ge\tau(r)+H\ge\tau(d)+H$, since $d\le r\Rightarrow\tau(d)\le\tau(r)$.
The register case is the same, with the adds below $\rho$.
\end{proof}

\begin{definition}[Holder]\label{def:holder}
Owner $p$ \emph{holds} $(e,r)$ at time $t$ iff its replica has $D(e)\ge r$.  The mirror is
never a holder.
\end{definition}

\begin{law}[Remove survival]\label{law:survival}
An unstable remove or register leaves a holder's replica only by
(i) $\ret_t$ at its stability,
(ii) an ack-gated move whose receiver holds it, or
(iii) a loss event.
No other program action drops it.  In particular, no eviction, cap or truncation does.
\end{law}

\begin{theorem}[No resurrection]\label{theorem:noresurrect}
Under Law~\ref{law:survival}, let $p$ hold $(e,r)$ at $t_0$, and let no loss event strike
$p$ before $\tau(r)+H+\sigma$.  Then no read at $p$ after $t_0$ serves an add covered by
$(e,r)$.
\end{theorem}
\begin{proof}
Before stability, $D_p(e)\ge r$ by Law~\ref{law:survival}, because no admissible transition
lowers it: a join takes the pointwise $\max$, and a move keeps it at the receiver.
Normalization then drops every covered add from each value $p$ stores or reads.  From
stability on, Lemma~\ref{lem:covers} retires every covered add, and every read applies
$\ret_t$.
\end{proof}

\begin{lemma}[Replication answers loss events, not eviction]\label{lem:replication}
A remove held by $k$ holders leaves the system before it is stable only if all $k$ drop it.
Under Law~\ref{law:survival}, that takes a loss event at each of them.  If eviction may drop
removes, one writer drops it from all $k$ with no loss event.
\end{lemma}
\begin{proof}
A writer may write to every placement of any topic, so it can fill the byte budget of each
of the $k$ owners.  Eviction is a function of the writes received.
\end{proof}

So \issue{920} breaks Law~\ref{law:survival} through a program action that a writer can
drive, and it breaks it at every replica at once.  The remaining sections discharge the law
and bound what it retains.

\section{Eviction and the factored carrier}\label{sec:evict}

\begin{definition}[Sides]\label{def:sides}
$\El(x)=(V,\emptyset,\bot,\beta)$ and $\Rm(x)=(\emptyset,D,\rho,0)$.
\end{definition}

The remove side carries the least stamped bound, $0$.  So a remove side read alone is live
only through $s(x)$, and it answers a lookup as a miss
(\texttt{bound\_live\_at} is false).  It never shadows a replica that holds elements.

\begin{lemma}[Factorization]\label{lem:factor}
For normalized $x$, $x=\El(x)\sqcup\Rm(x)$.  For every $E'$ whose adds are a subset of those
of some carrier the replica held and every $R'\sqsupseteq\Rm(x)$, $E'\sqcup R'$ contains no
add that $x$'s removes or register drop.
\end{lemma}
\begin{proof}
No add of $x$ is covered or below $\rho$, so normalization of the union drops nothing, and
$\max(\beta,0)=\beta$.  For the second claim, normalization of $E'\sqcup R'$ drops every add
that $R'$ covers or floors, and $R'$ covers or floors at least what $\Rm(x)$ does.
\end{proof}

\begin{law}[Eviction acts on elements]\label{law:evict}
Eviction maps $x\mapsto\Rm(x)$.  It may drop the element side of any carrier, and the whole
carrier iff $\Rm(\ret_t(x))=\bot$.
\end{law}

\begin{proposition}[What eviction may drop]\label{prop:maydrop}
Law~\ref{law:evict} satisfies Law~\ref{law:survival}, and it is the largest eviction that
does: dropping any unstable remove or register violates it.
\end{proposition}
\begin{proof}
$x\mapsto\Rm(x)$ keeps $D$ and $\rho$.  A remove that $\ret_t$ already dropped is stable.
Dropping an unstable $(e,r)$ is a program action outside (i)--(iii).
\end{proof}

\begin{law}[Authority per side]\label{law:authority}
A data carrier is stored as two records under one key.  $\El$ goes to the data store at
$(\Cached,\Rebuild)$, which is unchanged (Table~1 of~\cite{DurableRecords}).  $\Rm$ goes to a
\emph{removal store} at $(\Durable,\Retained)$, \texttt{RecordAuthority::Retained}.
\end{law}

\begin{proof}[Justification]
An absent remove side means ``never removed here, or retired at stability''.  That is the
meaning $\Retained$, whose obligations are N (no eviction) and D.  N is
Law~\ref{law:evict}.  D follows from Admissibility (Law~1 of~\cite{DurableRecords}).  D also
turns (ii) into a move under crashes: the receiver acknowledges a hand-off only after its
flush.  The element side keeps $\Rebuild$, because a lost add is an availability loss that
replication repairs.
\end{proof}

\begin{law}[Write order]\label{law:order}
$\mathrm{get}(k)=\El_k\sqcup\Rm_k$.  $\mathrm{put}(k,x)$ first writes $\Rm(x)$, or removes
the record if $\Rm(x)=\bot$, and does so only if it differs from the stored $\Rm_k$.  It then
writes $\El(x)$.  It returns $\Ok$ only after both, and an $\Err(\Exhausted)$ from the
removal store leaves both records unchanged.  $\mathrm{rm}(k)$ removes $\El_k$ and then
$\Rm_k$.
\end{law}

\begin{proposition}[Every crash state is safe]\label{prop:crash}
After any crash, $\mathrm{get}(k)$ serves no add that a remove or register the replica made
durable covers or floors, until that remove or register is stable.
\end{proposition}
\begin{proof}
D gives $\Rm_k\in\{\Rm(x_m),\Rm(x_{m+1})\}$ for the last two removal-store writes.  Under
$\Cached$, $\El_k$ is the element side of some earlier write, or $\bot$.  A remove side
changes only by joins, which grow it, and by $\ret_t$, which drops stable parts; so the
durable $\Rm_k$ covers or floors every unstable remove of every earlier carrier.  The claim
then follows from Lemma~\ref{lem:factor}, and Lemma~\ref{lem:covers} covers the stable parts.
A crash between the two removals of $\mathrm{rm}$ leaves a remove side alone, which is live
until stable and answers no lookup.
\end{proof}

\textbf{Cost.}  One flush per write that changes the remove side.  An $\mathrm{Extend}$ never
changes it, and after Section~\ref{sec:heartbeat} neither does a heartbeat.

\textbf{Budget.}  The removal store refuses and never evicts (Law~8 of~\cite{DurableRecords}).
Its budget $B_\Rm$ is shared by every topic at the owner, so it is partitioned by origin
(Law~\ref{law:share} of Section~\ref{sec:bound}).  An operation refused for room returns
$\Exhausted$ to its writer.  A refused sync join is never acknowledged, so the sender keeps its
copy and retries.

\textbf{Browsers.}  A browser's crash class is claimed, $\Durable^\circ$, and a store that
holds state for others is admissible at $(\Durable^\circ,\Retained)$ when $r\ge2$ holders keep
it and every repair period restores $r$ copies (Law~12 of~\cite{DurableRecords}, as amended by
\issue{1056}).  A remove side is held by every owner of its topic, $r=\mathtt{DATA\_REDUNDANT}=6$
by default, and repair restores the replica set.  So a browser opens the removal store in
IndexedDB at $(\Durable^\circ,\Retained)$, with strict commit and refusal.  By Proposition~9
of~\cite{DurableRecords}, a remove is lost only if every holder is lost within one repair
period.  A deployment configured with $r=1$ makes the browser removal store an inadmissible
point, as Law~12 makes any single-holder browser custody; native owners at
$(\Durable,\Retained)$ need no redundancy.

\textbf{IndexedDB that fails to open.}  Only then does a browser keep the remove side in
memory, refusing at its budget: $(\Volatile,\Retained)$, logged at \texttt{warn}.  Each restart
is the total loss of Proposition~2 of~\cite{DurableRecords}: the node comes back as a fresh
replica.  Theorem~\ref{theorem:noresurrect} then fails at that node alone, and by
Lemma~\ref{lem:replication} the remove survives at every other holder.

\textbf{Mirror.}  The fetch cache is $(\Volatile,\Rebuild)$.  Its count bound evicts whole
carriers, which ends a residency (Section~\ref{sec:cache}).  It drops no holder's remove, and
its read-repair delivers to a holder only values that the holder's removes still cover.

\section{The tombstone bound}\label{sec:bound}

\begin{definition}[Cap]\label{def:cap}
$\Tmax=\Dmax=\mathtt{ENTRY\_DATA\_MAX\_LEN}=1024$ removes per data carrier.
\end{definition}

\begin{law}[Remove admission]\label{law:admit}
(a) A peer-supplied delta $\delta$ with $|D_\delta|>\Tmax$ is not admissible.  Like the clock
bound, this is a predicate on the delta.
(b) A join or an operation whose result $y=\ret_t(x\sqcup\delta)$ has $|D_y|>\Tmax$ is refused
with $\Err(\Exhausted)$, and the slot is unchanged.
A data carrier's removes are never truncated.
\end{law}

\begin{proposition}[Refuse, not truncate]\label{prop:refuse}
Truncating a data carrier's removes violates Law~\ref{law:survival}.  A relay inbox may
truncate (\issue{913}), because it has one holder and one ack-gated relocation at a time.
\end{proposition}
\begin{proof}
Truncation drops an unstable remove by a program action, and Lemma~\ref{lem:replication}
applies, since a writer controls which removes exceed the cap.  A relay inbox has no second
replica that could keep a stale add.
\end{proof}

\begin{proposition}[One message]\label{prop:message}
Let $w_T$, $w_D$ and $w_\rho$ be the postcard suprema of one remove, one dot and one register
($147$, $115$ and $111$\,B: a version is a varint $u128$, a varint $u32$ counter and two
$42$-character DIDs, and a dot adds a varint
$u32$; a digest is $32$\,B).  Let $\lambda=3$ be the length prefix of a $32$\,KiB element and
$w_0$ the carrier's fixed fields.  Every admitted data carrier encodes in at most
\[
W=w_0+\Dmax(\Pmax+\lambda+w_D)+\Tmax w_T+w_\rho\approx33.83\,\text{MB},
\]
with $W_\Rm=w_0+\Tmax w_T+w_\rho$ for its remove side.  The carrier law is
$W+\mathtt{MAX\_PAYLOAD\_ENVELOPE\_OVERHEAD}\le\mathtt{TRANSPORT\_MAX\_SIZE}$
($60$\,MB), asserted at compile time.  Each $w$ is a constant witnessed with equality by a
test, as the envelope reserves are.
\end{proposition}
\begin{proof}
The cap $\Dmax$, the per-element bound $\Pmax$ and Law~\ref{law:admit} bound the counts, and
each field is at most its supremum.
\end{proof}

The margin is wide: removes cost $151$\,KB at the cap, against $26$\,MB of headroom.  We pick
$\Tmax=\Dmax$ for the size of a lookup reply, not for the message bound.

\textbf{Honest rate.}  The cap admits $\Tmax/(H+\sigma)\approx10$ removes per minute per topic,
sustained.  A topic that removes faster, such as a chat log that deletes each message, is
refused past the cap: its writer gets $\Exhausted$, and nothing is silently dropped.  Raising
$\Tmax$ would cost every lookup reply of a busy topic $w_T$ per remove, and every owner that
much of its removal budget.  An application that needs a higher rate shards the topic: $k$
topics $t_i$, with each element written to $t_{\dg(v)\bmod k}$, admit $k$ times the rate, and
since their keys differ they place at different owners.

\begin{lemma}[A sync delta simulates a writer]\label{lem:simulate}
Let $\delta$ be admissible at clock $c$.  Some writer, using at most
$2|D_\delta|+|V_\delta|+1$ operations stamped at clocks $\le c$, produces at an owner a
carrier that has the same coverage and remove count as $x\sqcup\delta$, up to dots in
$(c,c+\sigma]$.
\end{lemma}
\begin{proof}
Write each add with $\mathrm{Extend}$ and the register with $\mathrm{Overwrite}$.  Take a
remove $(e,r)$.  If $e=\dg(v)$ for an element $v$ that some carrier holds, then
$\mathrm{Tombstone}(v)$ covers what $(e,r)$ covers at dots up to $c$.  Otherwise no
held add has digest $e$ (collision resistance), so $(e,r)$ covers nothing, and the pair
$\mathrm{Extend}(v_e),\mathrm{Tombstone}(v_e)$ for a fresh $v_e$ adds one remove that covers
nothing held.  Admission bounds every $\tau(r)\le c+\sigma$.
\end{proof}

\begin{proposition}[Forgery through sync adds no power]\label{prop:forge}
A syncer cannot exceed $W$ (Law~\ref{law:admit}).  It covers an add only as a writer's
$\mathrm{Tombstone}$ would, with a dot up to $\sigma$ later.  It can fill one carrier's
removes, which refuses further removes on that carrier at that owner until
$c+\sigma+H+\sigma$.  By Lemma~\ref{lem:simulate}, any writer can do the same with
$O(\Tmax)$ operations.
\end{proposition}

So the defect of \issue{921} closes for one carrier.  A remote peer can still exhaust the
remove capacity of a topic, but every writer could already do that, since a data topic has no
write authority and its element cap already lets a writer displace the topic's elements.

\textbf{Isolation across topics.}  The removal store is one budget for every topic at an
owner, and it refuses instead of evicting, so junk held there stays until it is stable.  By
Lemma~\ref{lem:simulate}, one writer can place $\mathrm{Extend}(v_e),\mathrm{Tombstone}(v_e)$
pairs on fresh topics at an owner and, without a further bound, refuse every honest remove on
every topic there for $H+\sigma$.  A per-topic share does not prevent this: the junk topics
are fresh topics, each under its own cap.  The share must follow the sender.

\begin{definition}[Charge]\label{def:charge}
Owner $p$ charges each unstable remove $(t,e)$ it holds to an origin $o_p(t,e)$: the
authenticated sender of the transition that last raised $D_t(e)$ at $p$.  That is the
operation's writer, as the shell verified its signature, or the signer of the sync message.
A sync from a node that $p$'s own view names as an owner of $t$'s placement (the predecessor
handing off its range, or a replica peer) is charged to the \emph{holder class} $\star$
instead.  $o_p$ is local to $p$: it is stored beside each remove in the removal store, never
sent, and recounted at open.  Let $c_p(o)=|\{(t,e):o_p(t,e)=o\}|$ over unstable removes.
\end{definition}

\begin{law}[Origin share]\label{law:share}
A transition from sender $o\ne\star$ that would raise $c_p(o)$ above
$S=\Tmax$ is refused with $\Err(\Exhausted)$, and the slot is unchanged.  The store refuses
any transition past its total $B_\Rm=64\,\Tmax$ removes ($\approx9.6$\,MB at $w_T$).  A
transition that raises no $D_t(e)$ is never charged, and a charge is released when its remove
is retired at stability.
\end{law}

\begin{proposition}[Isolation]\label{prop:isolation}
Let an honest writer $w$ remove an element of topic $t$ at owner $p$.  The remove is refused
only if (i) $t$ holds $\Tmax$ unstable removes, (ii) $w$'s own share is exhausted,
$c_p(w)=S$, or (iii) the total is reached, which requires removes charged to at least
$\lceil(B_\Rm-c_p(\star))/S\rceil-1$ origins other than $w$, or removes charged to $\star$.
\end{proposition}
\begin{proof}
Laws~\ref{law:admit} and \ref{law:share} refuse for exactly three reasons: the topic cap, the
sender's share, and the total.  The total is
$c_p(\star)+\sum_{o\ne\star}c_p(o)$ with each $c_p(o)\le S$, so reaching it with
$c_p(w)<S$ needs the stated number of other origins.  An origin is the authenticated sender,
so no peer can charge $w$: a forged sync is charged to its own signer.
\end{proof}

Case (iii) needs either $63$ distinct identities or a node that $p$'s view makes an owner of
the topic.  The first is a Sybil attack and the second is a Byzantine storage owner, and
\texttt{SECURITY.md} lists both as non-goals.  The Sybil residual is the one of
\issue{907}, where free origins exhaust the replay table, and it closes there: when
\issue{907} prices or bounds origins, this share uses the same mechanism, not a second one.
Capacity exhaustion by one writer is thus
confined to its own share and to topics it fills to $\Tmax$; this is the availability residual
stated in \texttt{SECURITY.md}.  An honest writer may hold $S$ unstable removes per owner,
about $10$ per minute per owner, sustained.

\begin{proposition}[Refusal is safe]\label{prop:diverge}
Suppose owners $p$ and $q$ refuse each other's deltas because $|D_p\cup D_q|>\Tmax$.  Then
neither violates Theorem~\ref{theorem:noresurrect}.  They converge once the union fits, which
happens by $\max\tau(r)+H+\sigma$ over the removes present if no new remove arrives.
\end{proposition}
\begin{proof}
A refusal changes no slot, so no holder drops a remove.  An add that $p$ serves and $q$ has
removed is a remove $p$ never held.  That is a stale replica, not a resurrection.
\end{proof}

\section{The fetch cache}\label{sec:cache}

\begin{definition}[Cache machine]\label{def:cache}
The cache maps $c:K\rightharpoonup X$.  Its transitions run under one lock (the cache
transition), each at its own clock $t$:
\begin{align*}
\mathrm{put}(v)&:\ c[k]\leftarrow\ret_t\big(\live_t(c[k])\sqcup v\big),\\
\mathrm{obs}(o)&:\ c[k]\leftarrow\big(\live_t(c[k])\text{ or }\bot\big).o,\\
\mathrm{evict}(k')&:\ c[k']\leftarrow\bot,
\end{align*}
the last inside a put of a new key at the count bound.
A read at $t$ returns $\ret_t(c[k])$ when it is live and its bound has not elapsed, and it
changes nothing.
\end{definition}

\begin{law}[Pure reads]\label{law:pure}
A cache read is a function of one state and the reader's clock, and it never writes or
removes.  Every state change is one transition under the cache transition.  The store
orders single gets and puts atomically.
\end{law}

\begin{proposition}[Linearizable; no lost put]\label{prop:linear}
Each transition linearizes~\cite{HerlihyWing1990} at its store put, and each read at its store get.  After
$\mathrm{put}(v)$ at $k$, $c[k]\sqsupseteq\ret_t(v)$ holds until the next transition at $k$.
That transition is a put or obs, which keeps $c[k]$ ascending in $\sqsubseteq$ up to
retirement (\issue{982}), or the eviction of $k$, which ends the residency.
\end{proposition}
\begin{proof}
Transitions are totally ordered by the lock, and reads have no write effect.  So the history
is the lock-order sequence of transitions, with every read placed at its get.  The trace of
\issue{922} needs a read that stores, and no read stores.  The clocks need not be monotone in
lock order: by composition a state stored at an earlier clock is a less retired version of
the one at the later clock, and every read projects at its own clock.
\end{proof}

This is the state since \issue{913} made the cache read return its projection only.  A dead
cached value stays until a put replaces it or the count bound evicts it.  It costs one of
$\mathtt{LOCAL\_CACHE\_CAPACITY}$ slots and is served as absent.  Read-your-writes
(\issue{982}) holds per residency.  \issue{922} needs a test that witnesses
Proposition~\ref{prop:linear}, not a code change.

\section{Heartbeats}\label{sec:heartbeat}

Today a descriptor signs $\mathit{heartbeat\_at}$ and $\mathit{expires\_at}$, so its bytes
change at every heartbeat.  Each heartbeat therefore reads the registry, appends the new
descriptor and removes the old one.  With interval $I=30$\,s that leaves
$(H+\sigma)/I\approx200$ removes, about $25$\,KB, per registrant, and $N$ registrants read
about $N^2\cdot25$\,KB per interval.

\textbf{The issue's direction, stress-tested.}  \issue{918} proposes to drop the timestamps
and derive liveness from the element's dot.  A dot is stamped by whoever writes, and it is
unsigned: by Lemma~\ref{lem:simulate} a syncer can forge one too.  Any node that re-appends a
dead registrant's observed bytes would then refresh its dot, and so keep it listed forever.
Freshness that a peer cannot forge needs a signed bound.  The lease below keeps the bytes
stable and keeps that bound signed.

\begin{definition}[Lease descriptor]\label{def:lease}
A lease descriptor $d$ signs the registrant's identity fields, its process epoch, its
capabilities and a lease end $\ell$, and nothing that changes between heartbeats.  Its bytes
are fixed for the lease.  The lease length is $\Lambda=L=\mathtt{DEFAULT\_TTL\_MS}$ (10 min).
\end{definition}

\begin{law}[Heartbeat]\label{law:heartbeat}
At each heartbeat at clock $t$, a registrant writes exactly $\mathrm{Extend}([d])$ for its
current lease descriptor $d$.  It writes no remove and reads nothing.  If $t+2I\ge\ell$, or
if a signed field changed (restart, capabilities), it first signs $d'$ with $\ell'=t+\Lambda$
and appends $d'$ in place of $d$.  It never removes $d$.  A remove is written only on
withdrawal (advertising turned off), at most once per lease.
\end{law}

\begin{law}[Reader liveness]\label{law:reader}
Let $\tau^*(d)$ be the dot of $d$ in the fetched carrier.  A reader at clock $t$ lists, for
each DID, the descriptor with the greatest $\ell$ among those that verify and satisfy
\[
\ell\le t+\Lambda+\sigma,\quad T\le\Lambda,\quad t\le\ell,\quad t\le\tau^*(d)+T+\sigma,
\]
where $T$ is the descriptor's signed freshness window (the registry TTL, $90$\,s by default).
\end{law}

The first two conditions bound terms the signer chose: a lease longer than $\Lambda$, or a
window longer than the lease, is not a registry's, so no reader lists it.  A registrant
configured with $T>\Lambda$ is refused at start, since no reader would list it.

\begin{proposition}[Cost]\label{prop:hbcost}
Heartbeats write no removes.  A registrant writes at most one remove per withdrawal within
$H+\sigma$, so registries stay far below $\Tmax$.  A registrant leaves at most
$\lceil H/(\Lambda-2I)\rceil+1=13$ descriptors within the horizon.  The cap keeps the $\Dmax$
newest dots, and a heartbeat refreshes the current descriptor's dot, so superseded leases
are dropped first.  A heartbeat performs no registry read, which removes the $N^2$ read
traffic.
\end{proposition}
\begin{proof}
Law~\ref{law:heartbeat} writes $\mathrm{Extend}$ only, and an $\mathrm{Extend}$ of bytes
already held raises their dot (pointwise $\max$) without a remove.  A new lease is signed at
most once per $\Lambda-2I$.
\end{proof}

\begin{proposition}[Ghost bound]\label{prop:ghost}
Let a registrant stop at $t_s$.  If no node re-appends its bytes, it is listed until at most
$t_s+T+\sigma$.  Against any writer or syncer, each lease it signed is listed only at clocks
$t\in[\ell-\Lambda-\sigma,\ell]$, a window of at most $\Lambda+\sigma$, whatever $\ell$ the
signer chose; for an honest signer $\ell\le t_s+\Lambda$.
\end{proposition}
\begin{proof}
Its last honest dot has $\tau^*\le t_s$ and $T\le\Lambda$, which gives the first bound.  A
re-append or a forged dot can raise $\tau^*$, but not the signed $\ell$, and Law~\ref{law:reader}
requires $\ell-\Lambda-\sigma\le t\le\ell$, which gives the second.
\end{proof}

The cost is the trade-off of this law.  A replayed registrant can stay listed for up to
$\Lambda=10$\,min, where today it stays at most $T$.  In exchange, heartbeats produce no
removes, and the per-heartbeat registry read and its $N^2$ traffic disappear.  A larger
$\Lambda$ retains fewer descriptors and lengthens the ghost window.  The bound applies to
both registries (online nodes and onion exits).

\section{Implementation order}\label{sec:plan}

\begin{center}
\footnotesize
\begin{tabular}{@{}llp{0.6\columnwidth}@{}}
\toprule
PR & Issue & Law\\
\midrule
1 & \issue{922} & Law~\ref{law:pure}; a test witnessing Prop.~\ref{prop:linear}\\
2 & \issue{918} & Def.~\ref{def:lease}, Laws~\ref{law:heartbeat}, \ref{law:reader}
(\texttt{rings-node})\\
3 & \issue{921} & Def.~\ref{def:cap}, Law~\ref{law:admit}; the assertion of
Prop.~\ref{prop:message}\\
4 & \issue{920} & Laws~\ref{law:evict}--\ref{law:order}, \ref{law:share}; the removal
store\\
\bottomrule
\end{tabular}
\end{center}

PR~2 precedes PR~3.  Once PR~3 lands, today's heartbeats would hit $\Tmax$ with about five
registrants per topic, and their removes would be refused.  PR~4 needs PR~3, because the
removal store's budget is sized by $W_\Rm$.  PR~1 is independent.

\textbf{Operator steps}, in the CHANGELOG under \texttt{Unreleased}.  PR~2 changes the
descriptor encoding, so old descriptors no longer decode.  Readers skip them, and they
retire within $H$; this needs no step.  PR~4 adds the removal store (native directory and
browser database \texttt{dht-removals} beside the data store) and needs no step either.  A
carrier written before PR~4 keeps its removes in the data store until it is next written,
and the read join (Law~\ref{law:order}) still applies them.  An operator who wants
Law~\ref{law:evict} to hold from the first eviction removes the data store while the node is
stopped; replication rebuilds it.

\section{Summary of laws}

\begin{center}
\small
\begin{tabular}{@{}ll@{}}
\toprule
Question & Answer\\
\midrule
when stable & $t\ge\tau(r)+H+\sigma$ (Def.~\ref{def:stable})\\
what may drop it & $\ret_t$, a move, a loss event (Law~\ref{law:survival})\\
what eviction drops & the element side (Law~\ref{law:evict})\\
where removes live & removal store, $(\Durable,\Retained)$ (Law~\ref{law:authority})\\
how many & $\le\Tmax$, refused past it (Law~\ref{law:admit})\\
store isolation & per-origin share (Prop.~\ref{prop:isolation})\\
how large & $W\le$ one message (Prop.~\ref{prop:message})\\
forged by sync & simulates a writer (Lemma~\ref{lem:simulate})\\
cache order & pure reads, locked transitions (Law~\ref{law:pure})\\
heartbeat writes & $\mathrm{Extend}$ of a signed lease (Law~\ref{law:heartbeat})\\
\bottomrule
\end{tabular}
\end{center}

\sloppy
\bibliographystyle{IEEEtran}
\bibliography{removal-stability}

\end{document}
